\section{Experimentos}
Para realização de testes voltados para performance focando nos ORMs em diversos tipos de consultas e interações com a base de dados, baseou-se no parâmetro de tempo de processamento como medida de comparação. As conexões realizadas com a base de dados não representam diferenças expressivas, a grande diferença está no ORM e na hora de converter dados em objetos modelo mapeados juntamente com seus relacionamentos. Algumas decisões de desenvolvimento em relação aos ORMs podem torna-los custosos e lentos, por parte do Laravel destacam-se o uso de \textit{reflections}, citadas em \ref{rt_reflections}, que degradam a performance além de decisões pontuais como para buscar relacionamentos realizar duas consultas, uma para buscar os identificadores do modelo principal e outra para buscar os relacionamentos. As operações de conexão efetuadas com a base de dados não revelam disparidades significativas; entretanto, a distinção preponderante manifesta-se no desempenho dos ORMs e na fase de conversão de dados em objetos modelo, incluindo seus respectivos relacionamentos.